\documentclass[11pt,twocolumn]{article}
\usepackage[margin=1in]{geometry}
\usepackage{amsmath}
\usepackage{amssymb}
\usepackage{graphicx}
\usepackage{booktabs}
\usepackage{hyperref}
\usepackage{listings}
\usepackage{xcolor}
\usepackage{fancyhdr}
\usepackage{abstract}
\usepackage{titlesec}
\usepackage{cite}
\usepackage{microtype}

\hypersetup{
    colorlinks=true,
    linkcolor=blue,
    citecolor=blue,
    urlcolor=blue
}

\lstset{
    basicstyle=\ttfamily\footnotesize,
    backgroundcolor=\color{gray!10},
    frame=single,
    breaklines=true,
    language=Python
}

\pagestyle{fancy}
\fancyhf{}
\rhead{mycode, 2026}
\lhead{Personalized Local Coding Assistant}
\cfoot{\thepage}

\title{
    \textbf{mycode: A Personalized, Zero-Cost Local Coding Assistant} \\
    \large with RAG-Augmented Generation for Embedded ML Development
}

\author{
    Jordi Angel \\
    \textit{Owell, Ecuador} \\
    \texttt{github.com/jordiangel}
}

\date{May 2026}

\begin{document}

\maketitle

\begin{abstract}
The computational cost of frontier AI models continues to grow exponentially,
placing powerful coding assistance out of reach for resource-constrained
environments and offline deployments. I present \textbf{mycode}, a
terminal-based personal coding assistant that runs entirely locally at zero
ongoing cost, designed specifically for embedded machine learning development
at Owell, a startup operating in Ecuador. mycode combines DeepSeek-Coder-V2
16B, a Mixture-of-Experts architecture with only 2.4B active parameters, with
a Retrieval-Augmented Generation pipeline backed by PostgreSQL with pgvector,
a personal code snippet store, and a corpus of 34 domain-specific ML and
embedded systems papers totaling 2,612 embedded chunks. Empirical evaluation
on a custom C++17 benchmark covering algorithms, embedded systems, and machine
learning shows that RAG-based personalization improves pass rate from 60
percent to 80 percent over the bare model, closing roughly half the gap to
Claude Sonnet 4.6, which achieved 100 percent on the same benchmark. On a
project-scale task requiring a complete INT8 neural network inference engine
in C++17, mycode matched Sonnet 4.6 on compilation and code quality metrics,
with inference throughput of 100,000 forward passes in 2.75 milliseconds.
\end{abstract}

\section{Introduction}

The landscape of AI-assisted software development has been dominated by
cloud-hosted frontier models. GitHub Copilot, Claude Code, and OpenAI Codex
deliver high-quality code generation but impose ongoing subscription costs,
require reliable internet connectivity, and offer no mechanism for
personalization to an individual developer's patterns. For engineers working
in bandwidth-constrained or offline environments, such as embedded systems
development in Ecuador, these tools are often impractical.

Simultaneously, the compute required to train and serve frontier models
continues to scale aggressively. GPT-4 class models are estimated to require
thousands of A100 GPU-hours for a single inference pass at scale. The trend
toward ever-larger dense models exacerbates the gap between what frontier
labs can deploy and what a solo developer can run locally. Mixture-of-Experts
architectures have emerged as a promising counterpoint to this trend,
activating only a fraction of total parameters per forward pass and enabling
near-frontier performance at a fraction of the compute cost.

mycode was built to address a concrete engineering need at Owell, a startup
developing embedded ML systems in Ecuador. The system must function offline,
at zero marginal cost per query, and improve over time as the developer
accumulates domain-specific solutions. This paper describes the architecture,
implementation decisions, evaluation methodology, and benchmarking results of
mycode, including a direct comparison against Claude Sonnet 4.6, a
state-of-the-art frontier model.

\section{Background and Motivation}

\subsection{The Compute Scaling Problem}

Frontier model compute has followed a consistent scaling trajectory. Training
compute for leading models has roughly doubled every six to twelve months
since 2018. Inference costs scale proportionally with model size for dense
architectures. A 70B parameter dense model requires approximately 140GB of
GPU memory for full-precision inference, far exceeding the capacity of
consumer hardware.

For embedded ML development specifically, the mismatch is acute. Engineers
working on quantization, real-time inference, and hardware-constrained
deployment need coding assistance that understands INT8 arithmetic, fixed-point
representations, memory layout, and RTOS scheduling. Frontier models provide
this knowledge but at the cost of cloud dependency.

\subsection{Mixture-of-Experts as a Local Solution}

MoE architectures route each token through a small subset of expert networks
rather than the full model. DeepSeek-Coder-V2 16B activates only 2.4B
parameters per forward pass despite having 16B total parameters. This produces
inference throughput comparable to a 2-4B dense model while retaining the
representational capacity of a 16B model. On an Apple M3 Pro with 18GB
unified memory, this translates to approximately 67 tokens per second,
sufficient for interactive use.

\subsection{Retrieval-Augmented Generation}

RAG systems augment LLM generation by prepending retrieved context from an
external knowledge store to the model prompt. The retrieval step uses semantic
similarity between an embedded query vector and a database of pre-embedded
document chunks, enabling the model to reference domain-specific material not
present in its training data. For a personal coding assistant, this enables
two distinct use cases: retrieval of the developer's own past solutions, and
retrieval of relevant technical literature.

\section{System Architecture}

\subsection{Overview}

mycode consists of five layers: a CLI interface, an LLM inference layer, a
context gathering layer, a RAG pipeline, and a persistent knowledge store.
Figure 1 illustrates the data flow.

\begin{center}
\begin{minipage}{0.85\columnwidth}
\begin{lstlisting}[language=bash, caption=mycode query flow]
User query
    -> Context gathering (git diff, file)
    -> RAG retrieval (pgvector similarity)
    -> Conditional prompt injection
    -> DeepSeek-Coder-V2 16B (Ollama)
    -> Streamed response
    -> Quality gate validation
    -> Accept/reject save flow
    -> pgvector + Obsidian vault
\end{lstlisting}
\end{minipage}
\end{center}

\subsection{Model Layer}

The model is DeepSeek-Coder-V2-Lite-Instruct, a 16B MoE model with 2.4B
active parameters and 128K context length, served via Ollama on an Apple
MacBook Pro M3 Pro with 18GB unified memory. Measured performance on this
hardware shows a time-to-first-token of 99 milliseconds, a generation rate
of 67 tokens per second, and a prompt evaluation rate of 121 tokens per
second. Model cold load requires approximately 6 seconds. All subsequent
queries within the same session incur no reload overhead.

\subsection{Embedding and Vector Store}

Text embeddings are generated using \texttt{all-MiniLM-L6-v2} from the
sentence-transformers library, producing 384-dimensional vectors. This model
was selected for its low latency on CPU, requiring no GPU for embedding
inference. Embeddings are stored in PostgreSQL 16 with the pgvector extension,
using an IVFFlat index with cosine distance for approximate nearest neighbor
search. Two tables are maintained: \texttt{code\_snippets} for the developer's
accepted solutions, and \texttt{resources} for paper chunks.

\subsection{Knowledge Corpus}

The knowledge corpus consists of 34 domain-specific documents spanning
quantization methods, transformer architectures, embedded systems programming,
real-time operating systems, and linear algebra. Documents were converted from
PDF to markdown and segmented into 500-word chunks, yielding 2,612 total
chunks. Key documents include foundational papers on quantization,
specifically GPTQ, AWQ, LLM.int8, LoRA, QLoRA, and mixed-precision training,
as well as reference texts on FreeRTOS, real-time C++, and embedded systems
programming. All documents were selected for direct relevance to the embedded
ML development work at Owell.

\subsection{RAG Pipeline Design}

A critical design decision was made regarding when to inject retrieved context
into the model prompt versus presenting it as a reference panel to the user.
Early evaluation revealed that injecting paper chunks into the prompt
consistently degraded generation quality. The root cause was that prose
descriptions of technical concepts from papers, when injected as prior context,
caused the model to reproduce patterns from the retrieved text rather than
generating clean implementations. In one documented case, a retrieved code
snippet containing a bug (\texttt{int8\_t::max} instead of \texttt{INT8\_MAX})
was reproduced faithfully by the model in three consecutive generation attempts,
despite explicit instructions to the contrary.

The final pipeline implements a selective injection policy. Code snippets from
\texttt{code\_snippets} with cosine similarity above 0.55 are injected into
the prompt, as these represent the developer's own validated solutions to
similar problems. All paper chunk matches, regardless of similarity score, are
displayed to the user as a reference panel after generation but are never
injected into the model prompt. This preserves model generation quality while
surfacing relevant literature for the developer's reference.

\subsection{Quality Gate}

A compilation-based quality gate was added to the save flow after a buggy
snippet was found to have propagated through the pipeline. Before any code
snippet is persisted to the database or Obsidian vault, the system validates
it by syntax-checking Python via \texttt{py\_compile} or syntax-checking C++
via \texttt{g++ -fsyntax-only}. Code that fails validation is presented to
the user with the compiler error and requires explicit confirmation to save.
This prevents the knowledge base from accumulating incorrect patterns.

\subsection{CLI and Save Flow}

The CLI is implemented in Python using Click and Rich, invoked via a bash
wrapper at \texttt{/usr/local/bin/mycode}. After each generation, the user
is presented with an accept or reject prompt. Accepted solutions are embedded,
stored in \texttt{code\_snippets} with metadata including language, category,
and project, and saved as structured markdown to an Obsidian vault organized
in a hybrid topic and project hierarchy.

\section{Evaluation Methodology}

\subsection{Benchmark Design}

Standard coding benchmarks such as HumanEval and MBPP evaluate Python
function generation on general-purpose problems. These are insufficient for
evaluating a system designed for embedded ML development in C++. I designed
a custom evaluation suite with three components.

The first component is a 10-problem C++17 benchmark covering the three
primary domains of mycode usage: algorithms, embedded systems, and machine
learning. Problems include binary search, linked list reversal, merge sort,
circular buffer implementation, bitwise register manipulation, INT8
quantization, matrix multiplication, ReLU activation, stack from queues, and
Q16.16 fixed-point multiplication. Each problem is evaluated by compiling the
generated code against a hand-written test driver and checking for correct
output. The metric is pass at one, meaning a single generation attempt per
problem.

The second component is a project-scale evaluation requiring the model to
generate a complete, self-contained C++17 program implementing a lightweight
neural network inference engine for embedded deployment. Requirements include
INT8 weight storage, INT32 accumulation, ReLU and sigmoid activations, a
two-layer feedforward network with architecture Linear(2,4), ReLU, Linear(4,1),
Sigmoid, no heap allocation in the forward pass, and a benchmark measuring
100,000 forward pass calls. Scoring covers compilation, functional correctness
on XOR, code quality checks, and inference throughput.

The third component is a RAG quality benchmark measuring retrieval precision
across ten domain-specific queries, evaluating embedding score distributions,
pipeline latency, and the fraction of retrieved chunks scoring above the
injection threshold.

\subsection{Comparison Baseline}

Claude Sonnet 4.6 was used as the frontier model baseline. Sonnet 4.6 was
queried with identical problem specifications through the Claude web interface.
Generated code was saved locally and compiled using the same test drivers and
evaluation scripts used for mycode. This ensures a fair comparison on
functional correctness independent of any differences in output formatting.

\section{Results}

\subsection{Hardware Benchmarks}

Table 1 summarizes the inference performance of mycode on the M3 Pro.

\begin{table}[h]
\centering
\caption{Inference performance on Apple M3 Pro, 18GB.}
\begin{tabular}{lc}
\toprule
\textbf{Metric} & \textbf{Value} \\
\midrule
Model cold load & 6.0 seconds \\
Time to first token & 99 milliseconds \\
Prompt eval rate & 121 tokens per second \\
Generation rate & 67 tokens per second \\
RAG retrieval latency, avg & 194 milliseconds \\
RAG retrieval latency, min & 95 milliseconds \\
RAG retrieval latency, max & 428 milliseconds \\
\bottomrule
\end{tabular}
\end{table}

\subsection{RAG Retrieval Quality}

Table 2 summarizes embedding score distributions across ten domain-specific
queries.

\begin{table}[h]
\centering
\caption{RAG retrieval score distributions across ten queries.}
\begin{tabular}{lcc}
\toprule
\textbf{Metric} & \textbf{Resources} & \textbf{Code Snippets} \\
\midrule
Min similarity & 0.238 & 0.000 \\
Max similarity & 0.694 & 0.515 \\
Avg similarity & 0.506 & 0.161 \\
Above threshold & 26 of 30 & 5 of 30 \\
\bottomrule
\end{tabular}
\end{table}

The low average similarity for code snippets reflects the small size of the
snippet store at evaluation time, six total entries, rather than poor retrieval
quality. As the developer accumulates solutions through normal usage, this
distribution is expected to improve substantially.

\subsection{C++ Benchmark Results}

Table 3 presents pass rates on the 10-problem C++17 benchmark.

\begin{table}[h]
\centering
\caption{C++17 benchmark pass rates.}
\small
\begin{tabular}{p{1.4cm}ccc}
\toprule
\textbf{Domain} & \textbf{DS bare} & \textbf{mycode} & \textbf{Sonnet} \\
\midrule
Algorithms & 2/4 & 3/4 & 4/4 \\
Embedded   & 2/4 & 3/4 & 4/4 \\
ML         & 2/2 & 2/2 & 2/2 \\
\midrule
\textbf{Total}     & \textbf{6/10}  & \textbf{8/10}  & \textbf{10/10} \\
\textbf{Pass rate} & \textbf{60\%}  & \textbf{80\%}  & \textbf{100\%} \\
\midrule
Gen time   & 3.5s & 3.3s & ~2s \\
Cost/query & \$0  & \$0  & ~\$0.01 \\
\bottomrule
\end{tabular}
\end{table}

RAG personalization improved pass rate by 20 percentage points over the bare
model, from 60 percent to 80 percent. mycode with RAG closes approximately
half the gap to Sonnet 4.6. Average generation latency with RAG is marginally
lower than bare generation due to caching effects in the embedding model.

\subsection{Project-Scale Evaluation}

Table 4 presents results on the neural network inference engine project task.

\begin{table}[h]
\centering
\caption{Project-scale evaluation results.}
\begin{tabular}{lcc}
\toprule
\textbf{Metric} & \textbf{mycode} & \textbf{Sonnet 4.6} \\
\midrule
Compiled & Yes & Yes \\
XOR functional & No & Partial \\
Attempts required & 2 of 3 & 1 of 1 \\
INT8 weights & Yes & Yes \\
INT32 accumulation & Yes & Yes \\
No heap in forward & Yes & Yes \\
Activation functions & Yes & Yes \\
Benchmark present & Yes & Yes \\
Correct entry point & Yes & Yes \\
100k forward passes & 2.75 ms & 1.98 ms \\
Quality score & 68 of 100 & 78 of 100 \\
\bottomrule
\end{tabular}
\end{table}

Both systems compiled successfully and passed all six code quality checks.
Neither system achieved fully correct XOR classification. mycode produced
outputs of 0.5 for all four inputs, indicating the forward pass computation
collapsed to a constant. Sonnet 4.6 correctly classified two of four inputs
but produced 0.519 for inputs (0,0) and (1,1), which lie above the 0.5
decision boundary and were therefore classified incorrectly. On inference
throughput, mycode achieved 2.75 milliseconds for 100,000 forward passes
compared to 1.98 milliseconds for Sonnet 4.6, a difference of 0.77
milliseconds attributable to differences in generated implementation style
rather than hardware.

\section{Analysis}

\subsection{When RAG Helps and When It Hurts}

The evaluation revealed a consistent pattern. RAG improves performance on
targeted retrieval tasks where the developer has previously solved a similar
problem. On the binary search benchmark, the developer's saved solution scored
0.68 cosine similarity and was correctly injected, allowing the model to
reference a validated implementation. This pattern held across the algorithm
and embedded domains where prior solutions existed.

RAG consistently degraded performance on generative tasks requiring novel
composition, particularly when retrieved paper chunks were injected into the
prompt. The selective injection policy, restricting injection to code snippets
above 0.55 similarity and excluding all paper chunks from prompt injection,
resolved this degradation and produced the 80 percent pass rate.

\subsection{Quality Gate Impact}

The quality gate prevented one confirmed case of bug propagation. A saved
INT8 quantization snippet containing \texttt{int8\_t::max}, which is not valid
C++ syntax, was retrieved with 0.43 similarity on every quantization-related
query and caused the RAG-augmented model to reproduce the same compile error
across three consecutive generation attempts on the project task. Deleting the
buggy snippet and enabling the compilation-based quality gate eliminated this
failure mode entirely.

\subsection{Deployment Context}

mycode was designed for a specific operational context. Owell is developing
embedded ML systems in Ecuador, where internet connectivity may be unreliable
and cloud API costs accumulate at the scale of a small startup. Running
inference locally on an M3 Pro eliminates per-query costs entirely. The system
operates fully offline after initial model download. The learning loop, wherein
accepted solutions are embedded and stored for future retrieval, means the
system becomes more useful over time as domain-specific solutions accumulate.

\subsection{Limitations}

Several limitations of the current system are worth noting. The embedding
model, \texttt{all-MiniLM-L6-v2}, is a general-purpose model not fine-tuned
on technical or domain-specific text. Cosine similarity scores for technical
queries ranged from 0.24 to 0.69, lower than one might expect from a
domain-adapted model. Replacing it with a code-aware embedding model such as
\texttt{nomic-embed-text} may improve retrieval precision. The code snippet
store contained only six entries at evaluation time, limiting the impact of
the learning loop. With sustained usage over weeks or months, the retrieval
benefit is expected to grow substantially. Finally, the XOR project task
exposed a known limitation of locally quantized models on multi-step
compositional reasoning, where generating correct numerical constants requires
precise arithmetic that 2.4B active parameter models may not reliably perform.

\section{Related Work}

DeepSeek-Coder-V2 achieves 90.2 percent on HumanEval and 76.2 percent on
MBPP Plus for the full 236B variant, and 81.1 percent on HumanEval for the
16B Lite variant used in mycode. These scores were achieved on Python-centric
benchmarks. The custom C++17 benchmark used in this work is more representative
of embedded ML development tasks and shows a 60 percent bare pass rate for the
16B Lite model, consistent with the expectation that domain-specific
evaluations yield lower scores than general benchmarks.

RAG systems have been shown to reduce hallucination rates and improve factual
grounding in generation tasks. Research from Stanford indicates that poorly
evaluated RAG systems can produce hallucinations in up to 40 percent of
responses despite accessing correct information. The quality gate and selective
injection policy in mycode address this failure mode directly.

MoE architectures, as described in the DeepSeek-V2 technical report, offer
strong local performance within memory constraints relative to dense models of
comparable parameter count. The 16B MoE model with 2.4B active parameters
achieves generation rates competitive with dense 7B models while retaining
the knowledge capacity of a 16B model.

\section{Future Work}

Several directions are planned for continued development of mycode. First,
replacing \texttt{all-MiniLM-L6-v2} with a domain-adapted embedding model
such as \texttt{nomic-embed-text} is expected to improve retrieval precision
on technical queries and reduce the sensitivity of the similarity threshold.
Second, a query classifier to detect project-scale generative tasks and
automatically disable RAG injection would remove the need for manual
configuration. Third, as the Owell codebase grows, project-scoped retrieval
limiting search to solutions from the current project context will improve
relevance. Fourth, lightweight LoRA fine-tuning on the accumulated solution
corpus is a longer-term direction that could further personalize the model
beyond what RAG alone provides. Finally, integrating a continuous benchmark
harness that re-evaluates the system on the C++17 test suite after each
session would provide ongoing visibility into system quality as the knowledge
base grows.

\section{Conclusion}

I presented mycode, a personalized local coding assistant combining
DeepSeek-Coder-V2 16B MoE with a RAG pipeline backed by pgvector and a
domain-specific corpus of 34 embedded ML documents. Empirical evaluation on
a custom C++17 benchmark shows that RAG-based personalization improves pass
rate from 60 percent to 80 percent over the bare model, a 20 percentage point
improvement achieved at zero additional compute cost. On a project-scale
neural network inference engine task, mycode matched Claude Sonnet 4.6 on
compilation and code quality while running entirely offline at zero marginal
cost per query.

The key engineering insights from this work are three. First, RAG injection
must be selective: injecting paper chunks into the prompt degrades generation
quality, while injecting high-confidence matches from the developer's own
validated solutions improves it. Second, knowledge base quality matters more
than quantity: a single buggy snippet caused consistent downstream failures
across multiple queries, motivating the compilation-based quality gate. Third,
MoE architectures make locally competitive coding assistance genuinely
feasible on consumer hardware, closing roughly half the performance gap to
frontier models at a fraction of the operational cost.

This paper is available at \url{https://github.com/jangel19/jordiai}.

\end{document}
